\small
\begin{enumerate}[label=\textbf{\arabic*.}, leftmargin=0pt, itemsep=1.2em]
\item Что означает вещественное кодирование в генетическом алгоритме?
\begin{itemize}[label=$\square$, leftmargin=1.5cm, itemsep=0.2em, topsep=0.3em]
\item Особь хранится как вектор действительных чисел без бинарного кодирования
\item Все гены представляются только нулями и единицами
\item Каждая особь хранится как матрица ошибок
\item Алгоритм работает только с целыми числами
\end{itemize}
{\small\textit{Правильный ответ: Особь хранится как вектор действительных чисел без бинарного кодирования.}}

\item Какое преимущество вещественного кодирования по сравнению с бинарным особенно важно в непрерывной оптимизации?
\begin{itemize}[label=$\square$, leftmargin=1.5cm, itemsep=0.2em, topsep=0.3em]
\item Не требуется операция кодирования и декодирования переменных
\item Оно полностью исключает мутацию
\item Оно гарантирует точный глобальный минимум
\item Оно делает селекцию ненужной
\end{itemize}
{\small\textit{Правильный ответ: Не требуется операция кодирования и декодирования переменных.}}

\item В чем идея BLX-alpha кроссинговера?
\begin{itemize}[label=$\square$, leftmargin=1.5cm, itemsep=0.2em, topsep=0.3em]
\item Потомок выбирается из расширенного интервала между значениями родителей
\item Выбирается одна точка разреза битовой строки
\item Случайный сегмент разворачивается
\item Родители усредняются без разброса
\end{itemize}
{\small\textit{Правильный ответ: Потомок выбирается из расширенного интервала между значениями родителей.}}

\item Что определяет параметр alpha в BLX-alpha кроссинговере?
\begin{itemize}[label=$\square$, leftmargin=1.5cm, itemsep=0.2em, topsep=0.3em]
\item Насколько расширяется интервал выборки потомка за пределы родительских значений
\item Размер турнира в селекции
\item Число поколений до остановки
\item Число переменных в задаче
\end{itemize}
{\small\textit{Правильный ответ: Насколько расширяется интервал выборки потомка за пределы родительских значений.}}

\item Что делает гауссова мутация?
\begin{itemize}[label=$\square$, leftmargin=1.5cm, itemsep=0.2em, topsep=0.3em]
\item Добавляет к гену случайный шум с нормальным распределением
\item Меняет местами два случайных гена
\item Удаляет часть популяции
\item Всегда заменяет ген его средним значением
\end{itemize}
{\small\textit{Правильный ответ: Добавляет к гену случайный шум с нормальным распределением.}}

\item Зачем после гауссовой мутации часто выполняют ограничение значений переменных?
\begin{itemize}[label=$\square$, leftmargin=1.5cm, itemsep=0.2em, topsep=0.3em]
\item Чтобы особь оставалась внутри допустимой области поиска
\item Чтобы обнулить фитнес
\item Чтобы исключить кроссинговер
\item Чтобы преобразовать особь в бинарный код
\end{itemize}
{\small\textit{Правильный ответ: Чтобы особь оставалась внутри допустимой области поиска.}}

\item Как работает турнирная селекция в вещественном ГА?
\begin{itemize}[label=$\square$, leftmargin=1.5cm, itemsep=0.2em, topsep=0.3em]
\item Из нескольких случайно выбранных особей выбирается лучшая
\item Всегда выбирается первая особь массива
\item Выбираются только потомки предыдущего поколения
\item Каждая особь выбирается одинаково без учета качества
\end{itemize}
{\small\textit{Правильный ответ: Из нескольких случайно выбранных особей выбирается лучшая.}}

\item Что обеспечивает элитарная стратегия?
\begin{itemize}[label=$\square$, leftmargin=1.5cm, itemsep=0.2em, topsep=0.3em]
\item Сохранение лучших решений при переходе к следующему поколению
\item Полный отказ от мутации
\item Построение градиента функции
\item Преобразование вещественных особей в бинарные
\end{itemize}
{\small\textit{Правильный ответ: Сохранение лучших решений при переходе к следующему поколению.}}

\item Почему функция Растригина является трудной для оптимизации?
\begin{itemize}[label=$\square$, leftmargin=1.5cm, itemsep=0.2em, topsep=0.3em]
\item Она мультимодальна и содержит множество локальных минимумов
\item Она линейна и тривиальна
\item Она не имеет глобального минимума
\item Она определена только на бинарных значениях
\end{itemize}
{\small\textit{Правильный ответ: Она мультимодальна и содержит множество локальных минимумов.}}

\item Где расположен глобальный минимум функции Растригина?
\begin{itemize}[label=$\square$, leftmargin=1.5cm, itemsep=0.2em, topsep=0.3em]
\item В точке (0, 0)
\item В точке (1, 1)
\item На границе области поиска
\item У функции Растригина нет глобального минимума
\end{itemize}
{\small\textit{Правильный ответ: В точке (0, 0).}}

\item Какой эффект часто наблюдается при слишком малой вероятности мутации в мультимодальной задаче?
\begin{itemize}[label=$\square$, leftmargin=1.5cm, itemsep=0.2em, topsep=0.3em]
\item Популяция может преждевременно сойтись к локальному минимуму
\item Алгоритм перестает выполнять кроссинговер
\item Размер популяции автоматически растет
\item Значение alpha в BLX становится отрицательным
\end{itemize}
{\small\textit{Правильный ответ: Популяция может преждевременно сойтись к локальному минимуму.}}

\item Какой эффект возможен при слишком большой вероятности мутации?
\begin{itemize}[label=$\square$, leftmargin=1.5cm, itemsep=0.2em, topsep=0.3em]
\item Поиск становится слишком случайным и теряет хорошие решения
\item Сходимость всегда ускоряется
\item Все особи становятся одинаковыми
\item Глобальный минимум находится за одну эпоху
\end{itemize}
{\small\textit{Правильный ответ: Поиск становится слишком случайным и теряет хорошие решения.}}

\item Для чего нужна достаточно большая популяция в вещественном ГА?
\begin{itemize}[label=$\square$, leftmargin=1.5cm, itemsep=0.2em, topsep=0.3em]
\item Для лучшего покрытия непрерывного пространства поиска
\item Для устранения необходимости в фитнес-функции
\item Для полного исключения локальных минимумов
\item Для замены оператора кроссинговера
\end{itemize}
{\small\textit{Правильный ответ: Для лучшего покрытия непрерывного пространства поиска.}}

\item Что обычно сравнивают на графике сходимости генетического алгоритма?
\begin{itemize}[label=$\square$, leftmargin=1.5cm, itemsep=0.2em, topsep=0.3em]
\item Лучшее и среднее значение фитнеса по поколениям
\item Число входов и выходов сети
\item Только номера поколений
\item Размер файла результатов
\end{itemize}
{\small\textit{Правильный ответ: Лучшее и среднее значение фитнеса по поколениям.}}

\item Зачем строят траекторию лучшей особи на контурной карте функции?
\begin{itemize}[label=$\square$, leftmargin=1.5cm, itemsep=0.2em, topsep=0.3em]
\item Чтобы визуально увидеть движение поиска к минимуму
\item Чтобы заменить селекцию
\item Чтобы вычислить производную функции
\item Чтобы получить матрицу корреляций
\end{itemize}
{\small\textit{Правильный ответ: Чтобы визуально увидеть движение поиска к минимуму.}}

\item Какой ключевой компромисс должен обеспечивать BLX-alpha кроссинговер?
\begin{itemize}[label=$\square$, leftmargin=1.5cm, itemsep=0.2em, topsep=0.3em]
\item Компромисс между исследованием новых областей и использованием хороших родительских решений
\item Компромисс между бинарным и вещественным кодом
\item Компромисс между числом входов и выходов
\item Компромисс между классификацией и регрессией
\end{itemize}
{\small\textit{Правильный ответ: Компромисс между исследованием новых областей и использованием хороших родительских решений.}}

\item Почему вещественный ГА удобен для задач непрерывной оптимизации?
\begin{itemize}[label=$\square$, leftmargin=1.5cm, itemsep=0.2em, topsep=0.3em]
\item Потому что переменные представлены в естественной для задачи форме
\item Потому что он не использует случайность
\item Потому что он не требует фитнес-функции
\item Потому что он может работать только в одной размерности
\end{itemize}
{\small\textit{Правильный ответ: Потому что переменные представлены в естественной для задачи форме.}}

\item Что означает мультимодальность целевой функции?
\begin{itemize}[label=$\square$, leftmargin=1.5cm, itemsep=0.2em, topsep=0.3em]
\item Наличие множества локальных экстремумов
\item Наличие только одного экстремума
\item Линейный рост функции на всей области
\item Отсутствие глобального решения
\end{itemize}
{\small\textit{Правильный ответ: Наличие множества локальных экстремумов.}}

\item Какой из факторов помогает избежать стагнации популяции?
\begin{itemize}[label=$\square$, leftmargin=1.5cm, itemsep=0.2em, topsep=0.3em]
\item Поддержание разнообразия с помощью мутации и кроссинговера
\item Обнуление всех генов в конце поколения
\item Постоянное удаление лучших особей
\item Запрет случайности при селекции
\end{itemize}
{\small\textit{Правильный ответ: Поддержание разнообразия с помощью мутации и кроссинговера.}}

\item Какой критерий остановки часто используют в ГА?
\begin{itemize}[label=$\square$, leftmargin=1.5cm, itemsep=0.2em, topsep=0.3em]
\item Ограничение по числу поколений или отсутствие заметного улучшения
\item Только достижение нулевой ошибки классификации
\item Появление первого потомка
\item Обязательное совпадение всех особей
\end{itemize}
{\small\textit{Правильный ответ: Ограничение по числу поколений или отсутствие заметного улучшения.}}

\item Какой из операторов непосредственно создает новые комбинации генов на основе двух родителей?
\begin{itemize}[label=$\square$, leftmargin=1.5cm, itemsep=0.2em, topsep=0.3em]
\item Кроссинговер
\item Фитнес-функция
\item Редукция
\item Инициализация
\end{itemize}
{\small\textit{Правильный ответ: Кроссинговер.}}

\item Какой из операторов непосредственно добавляет локальные случайные возмущения решению?
\begin{itemize}[label=$\square$, leftmargin=1.5cm, itemsep=0.2em, topsep=0.3em]
\item Мутация
\item Элитарное копирование
\item Сортировка популяции
\item Нормализация данных
\end{itemize}
{\small\textit{Правильный ответ: Мутация.}}

\item Почему в вещественном ГА не требуется точность, определяемая длиной битовой строки?
\begin{itemize}[label=$\square$, leftmargin=1.5cm, itemsep=0.2em, topsep=0.3em]
\item Потому что переменные уже хранятся как вещественные значения
\item Потому что точность вообще не имеет значения
\item Потому что длина строки всегда равна 1
\item Потому что алгоритм работает только на целых числах
\end{itemize}
{\small\textit{Правильный ответ: Потому что переменные уже хранятся как вещественные значения.}}

\item Как обычно интерпретируется хорошая особь в задаче минимизации?
\begin{itemize}[label=$\square$, leftmargin=1.5cm, itemsep=0.2em, topsep=0.3em]
\item Это особь с меньшим значением целевой функции
\item Это особь с большей длиной хромосомы
\item Это особь с максимальным номером поколения
\item Это особь с наибольшим числом мутаций
\end{itemize}
{\small\textit{Правильный ответ: Это особь с меньшим значением целевой функции.}}

\item Почему эволюционные методы удобны для функций типа Растригина?
\begin{itemize}[label=$\square$, leftmargin=1.5cm, itemsep=0.2em, topsep=0.3em]
\item Они не требуют вычисления производных и способны исследовать сложный рельеф функции
\item Они работают только с выпуклыми функциями
\item Они полностью исключают локальные минимумы из функции
\item Они заменяют саму целевую функцию
\end{itemize}
{\small\textit{Правильный ответ: Они не требуют вычисления производных и способны исследовать сложный рельеф функции.}}

\item Что дает увеличение числа поколений при фиксированных остальных параметрах?
\begin{itemize}[label=$\square$, leftmargin=1.5cm, itemsep=0.2em, topsep=0.3em]
\item Больше возможностей для постепенного улучшения решений
\item Гарантированное нахождение глобального минимума за один шаг
\item Уменьшение размерности задачи
\item Отказ от кроссинговера
\end{itemize}
{\small\textit{Правильный ответ: Больше возможностей для постепенного улучшения решений.}}

\item Какой из параметров определяет интенсивность случайного изменения генов в гауссовой мутации?
\begin{itemize}[label=$\square$, leftmargin=1.5cm, itemsep=0.2em, topsep=0.3em]
\item Параметр sigma или связанный с ним коэффициент масштаба шума
\item Размер турнира
\item Число поколений
\item Количество особей в элите
\end{itemize}
{\small\textit{Правильный ответ: Параметр sigma или связанный с ним коэффициент масштаба шума.}}

\item Что наиболее характерно для хорошего набора параметров вещественного ГА?
\begin{itemize}[label=$\square$, leftmargin=1.5cm, itemsep=0.2em, topsep=0.3em]
\item Он обеспечивает баланс между сходимостью и сохранением разнообразия
\item Он исключает мутацию полностью
\item Он делает селекцию случайной
\item Он одинаков для любой задачи без изменений
\end{itemize}
{\small\textit{Правильный ответ: Он обеспечивает баланс между сходимостью и сохранением разнообразия.}}

\item Почему слишком маленькая популяция может быть проблемой?
\begin{itemize}[label=$\square$, leftmargin=1.5cm, itemsep=0.2em, topsep=0.3em]
\item Она хуже покрывает пространство поиска и легче вырождается
\item Она делает фитнес-функцию линейной
\item Она запрещает использовать BLX-alpha
\item Она автоматически увеличивает число локальных минимумов функции
\end{itemize}
{\small\textit{Правильный ответ: Она хуже покрывает пространство поиска и легче вырождается.}}

\item Какое утверждение о вещественном ГА наиболее корректно?
\begin{itemize}[label=$\square$, leftmargin=1.5cm, itemsep=0.2em, topsep=0.3em]
\item Это удобный инструмент глобальной оптимизации непрерывных функций со сложным рельефом
\item Он пригоден только для задач классификации
\item Он требует матрицы меток классов
\item Он не использует оператор селекции
\end{itemize}
{\small\textit{Правильный ответ: Это удобный инструмент глобальной оптимизации непрерывных функций со сложным рельефом.}}

\end{enumerate}
